% answers_short.tex —— 2015 年·春季上海卷 · 简短答案速查
% 由 tools/make_short_answers.py 自 answers.tex 的 \gksolhead/\gkanswer 提取，
% 并与 2015/2015_数学真题_春季上海_答案要点.md 逐题比对后生成，勿手改（改源 MD 后重跑）。

\begin{gknotice}
\textbf{使用说明}：本节只给最终答案便于快速核对，完整推导见解析版。答案为编写者独立推导，\textbf{非官方标准答案}。
\end{gknotice}

\gkpart{填空题}
\begin{itemize}[leftmargin=2.6em]
  \item[\textbf{1.}] \(\{3\}\)
  \item[\textbf{2.}] \(1-\mathrm{i}\)
  \item[\textbf{3.}] \(\pi\)
  \item[\textbf{4.}] \(\dfrac{1}{2}\)
  \item[\textbf{5.}] \((x-2)^{2}+(y-6)^{2}=1\)
  \item[\textbf{6.}] \(3\)
  \item[\textbf{7.}] \([3,4]\)
  \item[\textbf{8.}] \(2\)
  \item[\textbf{9.}] \(\{2\}\)
  \item[\textbf{10.}] \(84\)
  \item[\textbf{11.}] \(36\)
  \item[\textbf{12.}] \(y^{2}=2x-1\)
\end{itemize}

\gkpart{选择题}
\begin{center}\begin{tabular}{ccccccccccccc}
\toprule
题号 & 13 & 14 & 15 & 16 & 17 & 18 & 19 & 20 & 21 & 22 & 23 & 24 \\
\midrule
答案 & D & A & D & B & A & D & B & D & C & A & A & B \\
\bottomrule
\end{tabular}\end{center}

\gkpart{解答题}
\begin{itemize}[leftmargin=2.6em]
  \item[\textbf{25.}] \(8\)。
  \item[\textbf{26.}] \(a=0\)；最小值为 \(4\)，取得最小值时 \(x=2\)。
  \item[\textbf{27.}] 灯塔 \(B\) 与船相距约 \(4.2\) 海里；\(B\) 在船南偏东约 \(46^{\circ}\) 方向（即东偏南约 \(44^{\circ}\)，精确值 \(45.93^{\circ}\)）。
  \item[\textbf{28.}] (1) \(\dfrac{18}{5}\)；(2) \(b\in\left[\dfrac{\sqrt{22}}{2},3\right)\)
  \item[\textbf{29.}] (1) \(x<4\)；(2) \(k\leqslant \dfrac{25}{14}\)
\end{itemize}

\gkpart{附加题}
\begin{itemize}[leftmargin=2.6em]
  \item[\textbf{30.}] (C) 充要条件。
  \item[\textbf{31.}] (B) \([-4,0]\)。
  \item[\textbf{32.}] (D) \(\{a_{3n}\}\) 是等差数列。
  \item[\textbf{33.}] \(2\sqrt{2}\)。
  \item[\textbf{34.}] \(\sqrt{3}\)。
  \item[\textbf{35.}] $f(x)=\begin{cases}1, & 0<x<1,\\ x, & -1<x<0\end{cases}$（此时 $g(x)=\begin{cases}x, & 0<x<1,\\ -1, & -1<x<0\end{cases}$）。
  \item[\textbf{36.}] (1) 不存在；(2) \(x=\dfrac{\pi}{2}\)
\end{itemize}
